\chapter{Form før faget}

\begin{motto}
Ein snikkar treng ingen teori om tre.\\ Han treng tre, og hender som har halde tre lenge nok.
\end{motto}

Før ordet «design» fanst, fanst formverda likevel. Folk laga stolar ein kunne sitje på, krukker som heldt vatn, plogar som skar jord, hus som stod imot vinden, og dei laga dei utan ein einaste designar, utan eit einaste studium, utan ein einaste teori om form. Dei laga dei godt. I tusenvis av år vart verda forma kompetent av folk som ikkje hadde noko av det designfaget seinare kravde for å bli teke alvorleg, og som likevel produserte det meste av det vi i dag kallar god form. Forma kom før formgjevaren. Det er det ubehagelege fyrste faktumet, for det seier at faget ikkje var nødvendig for det faget seinare skulle hevde å eige.

Handverket les vi bakover, gjennom det som kom etter, og difor ser vi det feil. Vi tenkjer på det som det design ikkje lenger er, det manuelle, det lokale, det før-industrielle, eit ufullstendig design som venta på å bli ferdig. Men handverket venta ikkje på noko. Det var ein heil og fungerande kunnskapsmåte, ein som design seinare braut med og aldri klarte å erstatte med noko like haldbart. Glenn Adamson har vist at sjølve omgrepet «craft» er ein moderne oppfinning: ordet i si moderne tyding oppstår fyrst i kontrast til industri, altså fyrst etter at handverket alt var truga \autocite{adamson2013invention}. Namnet kom då tinget heldt på å forsvinne. Vi arvar det ordet, og med det ein synsvinkel som gjer handverkaren mindre enn han var.

\section{Tenking og laging i éin kropp}

Det avgjerande ved handverkaren er ikkje at han arbeidde med hendene. Det er at tenkinga og lagingane budde i den same personen, ofte i den same rørsla. Snikkaren som forma ein stol, formgav han ikkje fyrst og bygde han etterpå. Avgjerda om kor tjukk ein stolfot skulle vere, var ikkje teken på papir og overlevert til ei hand; ho var hand og avgjerd i eitt, teken i møte med dette stykket eik, denne åra, denne kvisten. Richard Sennett kallar dette einskapen av hovud og hand, og argumenterer for at det dugande ligg ikkje i at hovudet styrer handa, men i at handa tenkjer \autocite{sennett2008craftsman}. Kunnskapen sat ikkje over arbeidet som ein plan. Han sat i arbeidet som ei evne.

Difor var han òg dødeleg. Sennett peikar på verkstadene der ein kunnskap døydde med meisteren fordi han aldri vart skriven ned, slik den fullkomne fiolinen frå Cremona-verkstadene aldri vart laga om att etter at hendene som laga henne var borte: ingen klarte å rekonstruere det Stradivari visste, fordi det aldri låg nokon annan stad enn i hans eigne grep. Eit slikt tap er ikkje uhell. Det er den naturlege slutten på ein kunnskap som ikkje kan flyttast ut av kroppen som ber han. David Pye skilde mellom det han kalla risikoens handverk og visseheitens handverk, mellom arbeid der resultatet til kvar tid står på spel og avheng av utøvaren si dømekraft i augneblinken, og arbeid der resultatet er sikra på førehand av maskin eller mal \autocite{pye1968}. Pye sitt eige bilete er pennestreken mot trykket: handa som kan gå gale kvar einaste gong, mot maskina som ikkje kan gå gale i det heile. Handverkaren arbeidde under risiko, og difor kravde kvar rørsle merksemd, og denne merksemda var sjølv ein kunnskap som ikkje fanst nokon annan stad enn i utøvinga.

Kva slags kunnskap er dette? Ein som er fullstendig reell, snikkaren veit verkeleg noko om tre som den udanna ikkje veit, men som for det meste ikkje lèt seg formulere. Han kan ikkje seie deg kor mykje motstand han kjenner i høvelen før åra brest; han veit det berre, i armen, fordi han har høvla nok. Spør han kvifor han la fellinga akkurat slik, og han vil oftare svare «slik gjer ein det» enn gje deg ein grunn.

\section{Lauget som minne utan teori}

Korleis vart denne tause kunnskapen ført vidare, om han ikkje kunne skrivast ned? Ikkje gjennom bøker, men gjennom kroppar over tid. Lauget bar formkunnskapen i Europa frå mellomalderen og fram mot den industrielle revolusjonen, og det er det næraste handverket kom ein fagleg organisasjon. Overføringa gjekk gjennom lærlingordninga. Ein gut gjekk i lære hjå ein meister i mange år, fyrst som svein, til han kunne leggje fram eit meisterstykke og bli teken opp. Han lærte ikkje av lærebøker. Han lærte av å stå ved sida av meisteren, gjere det meisteren gjorde, få det retta, gjere det om att, til kunnskapen sat i hans eigne hender slik han sat i meisteren sine.

Dette produserte katedralar og kister og verktøy av ein kvalitet vi framleis beundrar, og det gjorde det utan ein felles, nedskriven teori om kvifor formene var som dei var. Lauget hadde reglar, strenge reglar om mål og materiale og kvalitet, men reglane var prosedyrar, ikkje forklaringar. Lauget visste \emph{korleis}, i ein grad få seinare fag har matcha; det visste ikkje, og trong ikkje vite, \emph{kvifor}. Adamson har synt korleis denne overføringsmåten ikkje var noko primitivt forstadium, men eit raffinert system for å halde taus kunnskap i live over generasjonar \autocite{adamson2007craft}. Lauget mangla ikkje noko det burde hatt. Det trong rett og slett aldri det fundamentet eit moderne fag treng, fordi det aldri kravde ein vitskaps autoritet. Det bad om eit verkstykke og tid, og leverte ein katedral.

\section{Det handa veit og maskina ikkje kan}

Då handverket fyrst kom under press frå maskina, fekk det sine fremste forsvararar, og dei sa skarpare enn nokon kva som budde i den tause kunnskapen. John Ruskin skreiv «The Nature of Gothic», sjette kapittel i andre bandet av \textit{The Stones of Venice}, at den gotiske bygningen er vakker nett fordi han ber spora av den ufullkomne, tenkjande, frie handa som laga han \autocite{ruskin1853stones}. Den jamne maskinforma er for Ruskin død, fordi ho ikkje ber minne om ein person som tok val; den ujamne handforma er levande, fordi kvar uregelmessigheit er ei avgjerd. Førti år seinare sette William Morris dette kapitlet på trykk for seg sjølv ved si Kelmscott Press, i 1892, og bygde eit heilt livsverk på innsikta, formulert som ein etikk: at arbeidet skal vere ei glede for den som utfører det, at skjønnleiken i tingen og verdigheita i arbeidet er to sider av same sak \autocite{morris1882hopes}. Den eine trykte den andre. Bandet mellom dei er ikkje ein parallell, men ei hand som rekkjer eit ark vidare.

Begge har auga for noko ekte. Det \emph{er} ein skilnad mellom forma som ber spor av ei dømande hand og forma som er stansa ut likt kvar gong, og skilnaden er ikkje innbilt. Men legg merke til kvar grunngjevinga deira hamnar. Når Ruskin skal seie kvifor den gotiske forma er betre, endar han ikkje i ein teori om form; han endar i ein moral om arbeid og ein metafysikk om liv. Når Morris skal grunngje den same dommen, endar han i ein politikk. Kunnskapen handa har, lèt seg ikkje, når han skal forsvarast i ord, formulere som kunnskap. Han kan berre omsetjast til overtyding. Det er ikkje fordi Ruskin og Morris var uklåre, dei var mellom dei klåraste i hundreåret, men fordi den tause kompetansen i sitt vesen ikkje har ein teoretisk botn å bli ført ned til. Krev du grunnen hans, får du ikkje ein grunn. Du får ei tru.

\section{Det vi veit og ikkje kan seie}

Det er ein innvending som ser ut til å redde fundamentet i siste liten: kanskje den tause kunnskapen \emph{kan} formaliserast, vi har berre ikkje gjort det enno; gjev oss tid og betre verktøy, og snikkaren si arm-kunnskap vil ein dag stå i ei lærebok. Mot dette står to argument som gjer ulikt arbeid, og som er verde å skilje, fordi dei altfor ofte vert stabla saman.

Det fyrste kjem frå Michael Polanyi, kjemikar i tjue år før han vende seg mot filosofien etter krigen, og som i Terry-førelesingane formulerte det som seinare stod att som ei setning: vi veit meir enn vi kan seie \autocite{polanyi1966tacit}. Poenget hans var ikkje at vi enno ikkje har sagt det, men at det ligg i kunnskapens natur at noko av han ber heilskapen utan sjølv å kunne nemnast, slik vi kjenner att eit andlet utan å kunne ramse opp dragen som utgjer det. Polanyi sitt eige skuledøme var sykkelen.

\begin{kasus}{Sykkelen ingen kan forklare}
Skriv ned, fullstendig og presist, korleis ein held balansen på ein sykkel. Det finst eit korrekt fysisk svar, og Polanyi gav det: for å rette opp ei helling skal styret dreiast slik at krumminga på banen er omvendt proporsjonal med kvadratet av farten, multiplisert med hellingsvinkelen. Ingen syklist på jord tenkjer dette, og ingen lærde det av ei slik setning. Du lærte det ved å falle, og falle, til armen visste. No har du kunnskapen fullkome i kroppen og kan ikkje gje han vidare i ord; den som vil lære, må sjølv opp på sykkelen og falle. Polanyi drog fram dømet for å vise nett poenget sitt: regelen finst, han er presis, og han er likegyldig, fordi det som ber prestasjonen ikkje er regelen, men det som ligg under han og aldri vart nemnt.
\end{kasus}

Det andre argumentet kjem frå Gilbert Ryle, som alt i 1949 skilde mellom \emph{knowing that} og \emph{knowing how}, å vite at noko er tilfellet og å vite korleis ein gjer noko \autocite{ryle1949mind}. Ryle hevda ikkje at kunnen ber ein unemneleg heilskap; han hevda at kunnen ikkje er anvend teori i det heile. Den dugande utøvaren handlar ikkje ved fyrst å slå opp ein regel i hovudet og så følgje han, for om han gjorde det, måtte han ha endå ein regel for korleis regelen skal brukast, og endå ein for den, i det uendelege. Kunnen er ikkje teori sett ut i livet. Ho er sin eigen ting. Der Polanyi seier at heilskapen ber utan å la seg nemne, seier Ryle at det ikkje er nokon teori der å nemne. Dei to drep den same draumen frå kvar si side: at handverkets \emph{knowing how} kan vekslast inn i eit fags \emph{knowing that} og leverast over disk til ein som aldri høvla, like lite som ein kan lære nokon å symje ved å overrekkje ei avhandling om hydrodynamikk. Harry Collins har seinare granska kor hard denne kjernen er, og synt at ein del av det tause er taust av prinsipp og ikkje av lettvinte grunnar, slik at det ikkje let seg flytte ut av kroppen som ber han \autocite{collins2010tacit}. Snikkaren si arm høyrer til den harde kjernen. Det er ikkje sikkert vekslinga er mogleg i det heile.

\section{Den dobbelte bokføringa}

Designfaget arvar handverkets kunnskapsform, den tause, blikk-baserte, etterliknande kompetansen frå atelier og verkstad, men det arvar henne saman med eit heilt anna sett krav. Det vil bli kalla ein disiplin. Det vil ha doktorgradar og professorat og tidsskrift, ein plass mellom dei lærde faga, retten til å tale med den autoriteten som følgjer av å vere eit fag og ikkje eit handlag. Og det er her bokføringa ikkje går opp. Du kan ikkje føre inntektene til ein vitskap og utgiftene til eit handverk i den same boka. Anten har du den tause, lokale kompetansen, og då gjer du krav som ein handverkar gjer, verdige krav, men avgrensa. Eller du har eit formalisert, delt, etterprøvbart fundament, og då kan du gjere krav som eit fag gjer. Det designfaget kjem til å prøve, gong på gong gjennom denne boka, er å krevje den eine stillinga på det andre grunnlaget. Det er ei dobbel bokføring, og som all dobbel bokføring fungerer ho berre så lenge ingen tel opp.

Eit ord om kva dette ikkje er. Det er ikkje ein dom over at design ikkje vart ein streng vitskap, for design er ikkje ein mislukka vitskap. Det er noko anna slag, og må grunngjevast på sine eigne premiss; å måle det mot ein moden vitskap og finne det for lett ville vere å gå glipp av poenget. Poenget er smalare og hardare. Den tause kunnskapen er ekte kunnskap, og verda hadde vore fattigare utan han; det er ikkje han som er problemet. Problemet oppstår fyrst når eit felt tek på seg den stillinga eit fag har, retten til å tale med autoritet og dele ut grader, og dermed pådreg seg pliktene eit fag har, og den fyrste av dei er å kunne gjere greie for kva det veit på ein måte som ikkje endar i «slik gjer ein det». Handverkaren var fri frå den plikta fordi han aldri tok på seg stillinga. Designfaget tek på seg stillinga og skyt frå seg plikta, og peikar på den tause kunnskapen som om det å \emph{ha} ein taus kompetanse var det same som å kunne gjere greie for han. Eit fundament er nett det som let ein gjere greie for ein taus kompetanse, slik anatomien lèt kirurgen seie kvifor handa gjorde det ho gjorde. Det er kva eit fundament \emph{er}, ikkje kva design måtte ha for å gjelde. Utan det blir påkallinga av taus kunnskap ei orsaking forkledd som eit forsvar, ein måte å sleppe å gje ein grunn ved å gje namnet til ein.

\vignett

Ein kunne innvende at faget aldri valde dette, at situasjonen vart påtvinga utanfrå. Det er sant, og det er nett det neste kapitlet handlar om. For så lenge tenking og laging budde i éin kropp, var det ingen som trong å spørje kva grunnlag forma kvilte på; grunnlaget var kroppen, og kroppen heldt. Spørsmålet om fundament oppstår fyrst i det augneblinket nokon riv tenkinga laus frå lagingane og gjev henne til ein annan person, i det augneblinket nokon teiknar det ein annan skal byggje. Det augneblinket har eit namn og ein dato, og det er der vi går no.

\laeringsmal{\begin{itemize}
\item Gjere greie for kvifor handverkets einskap av tenking og laging er ein eigen epistemisk struktur, og ikkje berre eit forstadium til moderne design.
\item Skilje mellom Pye sitt «risikoens» og «visseheitens» handverk og bruke skiljet til å plassere historiske produksjonsmåtar.
\item Skilje Polanyi sitt argument (heilskapen ber utan å la seg nemne) frå Ryle sitt (\emph{knowing how} er ikkje anvend teori), og forklare kva kvart av dei tyder for påstanden om at taus kunnskap kan formaliserast.
\item Identifisere figuren «den dobbelte bokføringa» og kjenne han att når han dukkar opp seinare i boka.
\end{itemize}}

\ovingar{\begin{enumerate}
\item Vel ein kvardagsleg dugleik du sjølv meistrar (sykle, knyte slips, kjenne når deig er ferdig elta). Prøv å skrive ned ein fullstendig instruksjon. Kva fell ut, og kvifor? Drøft kva dette seier om grensene for formalisering.
\item Ruskin grunngjev den gotiske forma med ein moral om arbeid, Morris med ein politikk. Finn ein moderne designtekst som grunngjev eit formval, og prøv å avgjere om grunngjevinga er teori eller moral i forkledning.
\item Lauget overførte kunnskap utan nedskriven teori og produserte likevel varig kvalitet. Diskuter: kva slags kunnskap \emph{krev} ikkje eit fundament? Er design eit slikt felt?
\item «Du kan ikkje føre inntektene til ein vitskap og utgiftene til eit handverk i den same boka.» Forsvar eller motbevis denne påstanden med eit konkret døme frå eit anna felt (sjukepleie, kokkekunst, musikk).
\end{enumerate}}

\vidarelesing{\fullcite{sennett2008craftsman}; \fullcite{pye1968}; \fullcite{adamson2007craft}; \fullcite{adamson2013invention}; \fullcite{polanyi1966tacit}; \fullcite{collins2010tacit}; \fullcite{ryle1949mind}; \fullcite{ruskin1853stones}; \fullcite{morris1882hopes}.}
